% エントリ 1(2026-11-03): 解析の設計
\section{[E1, 2026-11-03] 解析の設計}\label{sec:e01}

\subsection{問い}

GWTC-5 の BH スピンの集団は、四次の自己相互作用(強さ $\lambda=\mu^2/f^2$)をもつ超軽量スカラーの質量 $\mu$ と崩壊定数 $f$ を、どこまで排除するか。あるいは支持するか。特に、自由場では両立しない二つの観測事実を、ある範囲の $f$ が両立させるかを決める。
\begin{itemize}
\item \textbf{1 世代目の積み上がり}: 1 世代目の BH のスピンは、質量に依存する臨界値に集まっているように見える。
\item \textbf{2 世代目の速い回転}: 2 世代目らしい BH が $\chi\approx0.7$--$0.8$ で回っている。
\end{itemize}
テーマを選んだ経緯は \texttt{LVKtest/research/topic\_selection\_theory.md}(本文で照合したあとの「訂正」の節を含む)、決めた設計は \texttt{DESIGN.md} にある。

\subsection{現状(本文で照合)}

抜き書きは \texttt{Refs/notes/population\_papers.md} にある。

\paragraph{同じ GWTC-5 から逆の結論}
\begin{itemize}
\item \textbf{排除}(\arxiv{2607.01317}): 257 の合体で $1.7\times10^{-14}\le\mu\le3.3\times10^{-12}\eV$ を排除した(95\%)。
  \begin{itemize}
  \item 自然スピンは質量によらない単一の Beta 分布(同論文の式 2)。
  \item 超放射の率は水素原子近似の解析式で、$m=1$--6 を使う(End Matter の式 4)。
  \item 選択効果は入れていない。
  \item 自己相互作用は QCD アクシオンの $f_a$ でだけ確かめ、最終スピンが「場合により $\sim0.05$ 変わる」とした。ほかの $f_a$ は今後の課題としている(SM の「Superradiance for higher $\alpha$」)。
  \end{itemize}
\item \textbf{証拠}(\arxiv{2609.02678}): $\mu=0.71^{+0.18}_{-0.11}\times10^{-12}\eV$(90\%)、$\ln\mathcal{B}=7.8$(Table I)。
  \begin{itemize}
  \item 模型は、臨界スピン $\chi_\mathrm{SR}(M_f,T_\mathrm{age},\mu)$ を超える確率を、$\chi_\mathrm{SR}$ を中心とする幅 $\sigma_\mathrm{SR}$ の Gauss 分布に寄せる(同論文の式 1)。$\sigma_\mathrm{SR}=0.18^{+0.05}_{-0.03}$。
  \item $\chi_\mathrm{SR}$ は SuperRad で計算する。年齢 $T_\mathrm{age}$ は事象ごとに $[1\,\mathrm{Myr},2.2\,\mathrm{Gyr}]$ で対数一様。
  \item 選択効果は LVK の注入で入れている。
  \item 比べた帰無模型は LVK の Gauss 成分のスピン模型だけで、階層合体の模型はない。自己相互作用は入れていない。
  \end{itemize}
\end{itemize}

\paragraph{分かれる原因}
両者とも自由場である。分かれる原因は積み上がりの幅と、一つの事象である。
\begin{itemize}
\item \arxiv{2609.02678} は自分の模型の幅を 0 にした場合(デルタ関数)も試している。その場合 $3.5\times10^{-14}$--$1.6\times10^{-12}\eV$ がほぼ排除され、それは GW241011 だけによると書いている(End Matter III)。
\item GW241011 の主星は $m_1=19.6^{+3.6}_{-2.5}\Msun$、$\chi_1=0.78\pm0.09$ で、軌道から傾いている。LVK はこれを前の合体の残骸(2 世代目)と解釈している(\review\ T04-28)。
\end{itemize}

\paragraph{空いているところ}
\begin{itemize}
\item LVK のスピンの集団を $(\mu,f)$ で走査した解析はない(2025-01 以降を INSPIRE の被引用検索で確認)。最新のレビュー \arxiv{2610.11091} の \S V も、集団解析は自由場のままだと書いている。
\item 2 世代目の BH の自然スピンを超放射の較正に使った解析もない。
\item 近い仕事:
  \begin{itemize}
  \item \arxiv{2609.12052}: GWTC-5 で $(\mu,1/f)$ 面を扱う。ただし観測量はスピンではなく、銀河核での $q<0.3$ の合体率。
  \item 個々の事象による $(\mu,f)$ の制限: GW231123 と GW190517(\arxiv{2507.20979}、\arxiv{2507.21788})、GW190412 と GW190517(\arxiv{2503.10347})。
  \item 自由場での一つの事象による制限: GW241011(\arxiv{2510.26931})、O4b・O4c の速く回る事象(\arxiv{2608.11620})。
  \end{itemize}
\end{itemize}

\subsection{仮説の形}

自由場では、$\mu\approx7\times10^{-13}\eV$ の雲は、$15\Msun$、$\chi=0.69$ の BH の回転を約 300 年で臨界値まで奪う(E2)。2 世代目の BH が生まれてから次に合体するまでの時間(遅れ)は Myr 以上なので、GW241011 は残らない。

自己相互作用があると、雲は準位間のエネルギー交換と無限遠への放出で準平衡に達し、回転を失う速さは $f^2$ に比例して遅くなる(\arxiv{2011.11646})。すると次のような $f$ の範囲がありうる。
\begin{itemize}
\item 年をとった 1 世代目の BH(遅れは Gyr 程度)は回転を失う。
\item 若い 2 世代目の BH(遅れは 10--100 Myr 程度)は回転を保つ。
\end{itemize}
回転を失う途中の BH が年齢の分布に応じて散らばるなら、\arxiv{2609.02678} が手で入れた幅 $\sigma_\mathrm{SR}$ は物理的に説明できる。この範囲がなければ、その質量はどの $f$ でも排除される。どちらの結論も論文になる。

\subsection{理論の模型}

自己相互作用のある多準位の雲の運動論を解き、最終スピンの写像 $\chi_f(M,\chi_i,\mu,f,t)$ を表にする。
\begin{itemize}
\item \textbf{準位}: まず $211$、$322$、$433$ と、交換で生まれる減衰する準位(\arxiv{2011.11646} の組)。次に $n\le5$ に広げる(\arxiv{2412.03655})。
\item \textbf{過程}: 超放射、準位間の交換(Bose 増強 $\propto N_iN_j(1+N_k)$)、BH への吸収、無限遠へのスカラーの放出、重力波の放出(消滅と準位間の遷移)。
\item \textbf{率}: 超放射の率は、小さい $\alpha$ では解析式(E2)、$\alpha\gtrsim0.2$ では相対論的な率を使う。相対論的な率は SuperRad の当てはめ(\arxiv{2211.03845}、BSD-3)か、\arxiv{2412.03655} の公開コード(Julia、ライセンスの表示なし)と照合する。どちらを使うかはライセンスを確かめてから決める。
\item \textbf{解き方}: 準位間の交換は超放射よりずっと速く、方程式は硬い。占有数の対数を変数にして陰的な解法を使い、準平衡の解析解(\arxiv{2011.11646})と照合する。
\end{itemize}

\subsection{集団の模型}

\begin{itemize}
\item \textbf{二つの世代}: 1 世代目と 2 世代目の混合とする。
  \begin{itemize}
  \item 1 世代目の自然スピンは柔らかい模型にする(Beta 分布と、質量に依存する混合)。
  \item 2 世代目の自然スピンは、数値相対論の残骸の公式から、親の質量比とスピンの分布で決める($\approx0.69$ の周り)。
  \end{itemize}
\item \textbf{年齢}: 1 世代目は星形成から合体までの遅れ、2 世代目は前の合体から次の合体までの遅れとする。それぞれ星団模型の範囲で周辺化する。
\item \textbf{超放射の写像}: 生まれたときのスピンと年齢を、理論の写像で観測時のスピンに移す。
\item \textbf{選択効果と質量}: 選択効果は GWTC-5 の注入で入れる。質量・赤方偏移の模型と同時に推定する。
\item \textbf{比べる模型}: 次の四つを比べる。
  \begin{itemize}
  \item ボソンなし(柔らかい帰無模型。2 世代目を含む)
  \item 自由場
  \item 自己相互作用あり
  \item \arxiv{2609.02678} の模型の再現(このノートの部品で組む)
  \end{itemize}
\end{itemize}

\subsection{データ}

GWTC-5 の事後分布と累積の注入を使う(\texttt{data/README.md})。PBH\_LVKpop と同じファイルなので、置き場所は共用にする方向でユーザーに確かめる。\arxiv{2609.02678} の事後分布(Zenodo 22242056、146 MB)は、照合のために取得を検討する。

\subsection{照合の計画}

\begin{itemize}
\item E2: 自由場の超放射の率を、小さい $\alpha$ の極限と文献の値と照合する。半解析の飽和の扱いを ODE と照合する。
\item E3: 自己相互作用のある運動論。\arxiv{2011.11646} の準平衡の占有数と、回転を失うかどうかの $f$ の境界を再現する。
\item E4: 自由場の集団解析で、\arxiv{2607.01317} の排除と \arxiv{2609.02678} の証拠の両方を、このノートの部品で再現する。違いが積み上がりの幅と GW241011 に帰着することを確かめる。
\end{itemize}

\subsection{完成の条件}

\begin{enumerate}
\item $(\mu,f)$ 面の排除または証拠の地図。QCD アクシオンの線と、ALP として動機のある $f\sim10^{14}$--$10^{17}\GeV$ を描く。
\item \arxiv{2609.02678} の質量が、ある $f$ で生き残るかの判定。生き残るなら、その $f$ の範囲と、予言される重力波の信号。
\item 頑健さ: 遅れの分布、1 世代目の自然スピン模型、2 世代目の割合、波形の系統誤差、GW241011 と GW231123 を入れるか除くか。
\item 自由場の場合に、先行研究の二つの結果をこのノートの部品で再現した比較。
\end{enumerate}

\subsection{日程}

\begin{itemize}
\item 11 月前半: E2(自由場)、E3(自己相互作用の運動論)、写像の表。
\item 11 月後半: データの取得、E4(自由場の集団解析の再現)。
\item 12 月: 自己相互作用ありの集団解析、頑健さ。
\item 2027 年 1 月: 原稿。
\end{itemize}

\begin{caveat}
\begin{itemize}
\item 2 世代目の遅れの下限が、較正の強さを決める。\arxiv{2107.11730} は、2 世代目の多くは若すぎて雲が育たないとした。ただしそれは $\mu=4.1\times10^{-14}\eV$ での話である(同論文 \S4.3)。$\mu\sim10^{-12}\eV$ の $10$--$20\Msun$ では雲は数百年で育つ(E2)。
\item 自己相互作用で回転を失う速さが遅くなる範囲では、雲の放出する重力波とスカラー波も変わる。連続波の探索の制限と矛盾しないかを、結果の $f$ の範囲で確かめる。
\item 先を越される危険がある。2--3 か月で出す。
\end{itemize}
\end{caveat}
